\section{Integración Sin Código: El Plugin \texttt{dashai-frankenstein}}
\label{sec:dashai-plugin}

El CLI basado en esquemas y el generador web descritos en la Sección~\ref{sec:system-design} sirven a investigadores cómodos con archivos de configuración YAML. Una audiencia complementaria---especialistas de dominio, educadores y analistas que necesitan un modelo entrenado pero no un script de entrenamiento---está mejor servida por un entorno gráfico sin código. Para alcanzar a esa audiencia, Frankenstein Transformer distribuye un paquete adaptador complementario, \texttt{dashai-frankenstein}, que registra las clases de modelo de Frankenstein como componentes de \textbf{dashAI} (\texttt{\url{https://github.com/DashAISoftware/dashAI}}), una aplicación de escritorio open-source y local-first para aprendizaje automático desarrollada por un consorcio académico chileno (Universidad de Chile/FCFM, CENIA, IMFD y financiamiento ANID). Como ambos sistemas son basados en esquemas en esencia---dashAI genera sus formularios de configuración directamente desde esquemas Pydantic, mientras que Frankenstein valida su YAML contra un JSON Schema estricto---la integración no requirió cambios en el backend de dashAI: el plugin se descubre mediante el grupo de entry-points estándar \texttt{dashai.plugins} al inicio, exactamente como documenta la guía de plugins de dashAI (\texttt{\url{https://docs.dash-ai.com/learn/guides/plugins/}}), donde los plugins son paquetes PyPI ordinarios cuyas capacidades generativas de primera parte se distribuyen ellos mismos como plugins.

\subsection{Componentes Registrados}

Un único grupo de entry-points \texttt{dashai.plugins} registra cinco componentes en el registro de componentes de dashAI, cubriendo toda la amplitud de las clases de modelo del toolkit:

\begin{itemize}[leftmargin=*]
    \item \textbf{\texttt{FrankensteinMLMModel}} (\texttt{BaseModel}) --- se vincula a la \texttt{TextClassificationTask} de dashAI. Entrena el encoder de Frankenstein con una cabeza de clasificación de precisión completa con pooling sobre la salida del encoder (Estrategia~A de la auditoría de integración), heredando todos los mezcladores de secuencia, normalizadores y familias de optimizadores del motor subyacente; \texttt{predict} devuelve una matriz de probabilidades por clase.
    \item \textbf{\texttt{FrankensteinDecoderModel}} (\texttt{BaseGenerativeModel}) --- se vincula a \texttt{TextToTextGenerationTask}. Fuerza \texttt{model\_class: frankensteindecoder} (que a su vez fuerza \texttt{mode: decoder}) y expone muestreo autorregresivo top-\emph{k}; los parámetros de generación (\texttt{max\_new\_tokens}, \texttt{temperature}, \texttt{top\_k}) permanecen como campos nativos del formulario de dashAI porque son ajustes de tiempo de inferencia sin lugar en el esquema de entrenamiento de Frankenstein.
    \item \textbf{\texttt{FrankensteinViTClassifier}} (\texttt{BaseModel}) --- se vincula a \texttt{ImageClassificationTask}, reutilizando la cabeza de clasificación integrada del Vision Transformer (Apéndice~\ref{sec:annex-vision}).
    \item \textbf{\texttt{FrankensteinViTSegmenter}} (\texttt{BaseModel}) --- se vincula a una \texttt{SegmentationTask} proporcionada por el plugin, una nueva componente \texttt{BaseTask} que dashAI no incluye; usa la cabeza de segmentación del ViT y devuelve mapas de clases por píxel.
    \item \textbf{\texttt{SegmentationTask}} (\texttt{BaseTask}) --- la tarea complementaria que declara tipos de columna imagen-entrada/máscara-salida para el segmentador.
\end{itemize}

Los cinco componentes comparten una fachada delgada sobre la API de motor reutilizable de Frankenstein (la interfaz sin CLI y sin supervisor extraída en la Fase~0 del plan de integración), de modo que existe exactamente una dependencia pesada (\texttt{frankenstein-transformer}) detrás del plugin.

\subsection{Esquema de Configuración de Passthrough}

La integración preserva la filosofía basada en esquemas de Frankenstein (Sección~\ref{sec:system-design}) manteniendo el JSON Schema de Frankenstein como la única fuente de verdad en lugar de traducir sus más de $150$ campos de modelo y $38$ campos de entrenamiento a campos nativos de formulario de dashAI---una traducción que sería frágil frente a un esquema en evolución activa. En su lugar, cada componente expone un esquema Pydantic mínimo cuyo único campo visible para el usuario, \texttt{frankenstein\_json}, es una cadena que contiene una configuración completa de entrenamiento de Frankenstein como documento \emph{JSON de una sola línea} (una concesión a los campos de texto de una sola línea de dashAI; los usuarios construyen el YAML con el generador web basado en esquemas y lo convierten con una sola línea). Antes de que cualquier entrenamiento o construcción de modelo se lance, el adaptador ejecuta un pipeline de validación de tres etapas:

\begin{enumerate}[leftmargin=*]
    \item \textbf{Parseo JSON} --- detecta errores de sintaxis en la carga de una sola línea.
    \item \textbf{Validación JSON Schema} --- \texttt{jsonschema.validate} contra el esquema de Frankenstein resuelto, imponiendo \texttt{additionalProperties: false} y todos los rangos de enums (mezcladores, optimizadores, normalizaciones, activaciones).
    \item \textbf{Validación del cargador de configuración} --- el propio \texttt{load\_training\_config} de Frankenstein impone las restricciones cruzadas entre componentes (divisibilidad de \texttt{hidden\_size} entre \texttt{num\_heads}, \texttt{num\_kv\_heads} dividiendo \texttt{num\_heads}, acuerdo vocabulario/tokenizador, reglas de flags BitNet, presencia del optimizador por tarea, \texttt{frankensteindecoder} forzando \texttt{mode: decoder}).
\end{enumerate}

Cualquier fallo se presenta al usuario de dashAI como un único \texttt{ValueError} legible etiquetado con la etapa fallida, de modo que las configuraciones inválidas nunca alcanzan el bucle de entrenamiento---el mismo contrato fail-fast que provee el CLI. La inyección de datasets la manejan adaptadores que mapean el \texttt{DashAIDataset} de dashAI (un envoltorio de Hugging Face \texttt{datasets}) sobre las expectativas de dataset de Frankenstein, inyectan \texttt{vocab\_size} desde el tokenizador resuelto para que la matriz de embeddings coincida con el vocabulario, y transmiten las métricas por época de vuelta al almacén de métricas de dashAI para que el progreso del entrenamiento sea visible en la UI.

\subsection{Trabajo en Proceso}

El plugin es un trabajo en proceso. Sus fases están implementadas---el MVP de clasificación de texto MLM (Fase~1), el decoder generativo y el clasificador ViT (Fase~2), y el segmentador ViT más la \texttt{SegmentationTask} proporcionada por el plugin (Fase~3)---pero el paquete aún no se ha ejecutado dentro de una instancia viva de dashAI: el contrato de componentes se implementó a partir de una auditoría del registro de plugins, las clases base y la arquitectura de cola de trabajos de dashAI en lugar de validarse de extremo a extremo. La publicación en PyPI también está pendiente, tanto del propio plugin como de su piso de dependencia (\texttt{frankenstein-transformer>=1.1.0}, el release con la API de motor que reemplaza al \texttt{1.0.0} pre-motor ya en PyPI). Las limitaciones conocidas que se arrastran a la siguiente iteración incluyen: el componente MLM ejecuta un bucle de clasificación supervisada directa en lugar de una receta de preentrenamiento MLM seguido de ajuste fino; el esquema de passthrough difiere los campos nativos curados para ajustes de alto impacto (clase de modelo, patrón de mezclador, familia de optimizador) a una versión futura; y la capacidad de dashAI para renderizar columnas de salida de máscaras de segmentación no está confirmada y puede requerir una adición local al frontend. Estos ítems están registrados en la auditoría de integración del repositorio (\texttt{\url{https://github.com/erickfmm/frankenstein-transformer}}), y el diseño deja intencionalmente el núcleo del backend de dashAI intacto, de modo que futuros releases de dashAI no puedan romper el adaptador mediante acoplamiento.

La integración completa la historia de accesibilidad del toolkit: el mismo contrato de configuración validado que impulsa el CLI, el generador Streamlit y el pipeline de exportación ahora también impulsa una aplicación de trabajo sin código de terceros, permitiendo que el espacio arquitectónico de Frankenstein se explore desde una interfaz gráfica sin sacrificar las garantías de esquema de las que depende la reproducibilidad.